%------------------------------------------------------------------------------%
\documentclass[fleqn,utf8,aspectratio=169,12pt]{beamer}

\usepackage{calculo_varias_variaveis-1}

\title{Inexistência do Limite de Funções de Várias Variáveis}

%------------------------------------------------------------------------------%
\begin{document}

\addtitle

%------------------------------------------------------------------------------%
\section{Motivação}

%------------------------------------------------------------------------------%
\begin{question}
  Analise o limite de \(f(x, y) = \frac{y}{x}\)
  quando \((x, y)\)  se aproxima de \((0, 0)\)

  \pause
  \vspace{\baselineskip}
  A reta \(x=0\) não faz parte do domínio de $f$

  \pause
  O ponto \((0, 0)\) não está no domínio de $f$

  \pause
  Não podemos usar as propriedades diretamente

  \pause
  Não existe manipulação que remova a indeterminação
\end{question}

%------------------------------------------------------------------------------%
\begin{resolution}{Solução}
  Analisando $f$ sobre a reta \(y=x\)
  \pause
  \[
    f(x, y) \evalat{y=x}{}
    \pause
    = f(x, x)
    \pause
    = \frac{x}{x}
    \pause
    = 1
    \pause
    \qquad\text{ pois } x=0 \text{ não está no domínio de } f
  \]

  \pause
  Analisando $f$ sobre a reta \(y=0\)
  \pause
  \[
    f(x, y) \evalat{y=0}{}
    \pause
    = f(x, 0)
    \pause
    = \frac{0}{x}
    \pause
    = 0
    \pause
    \qquad\text{ pois } x=0 \text{ não está no domínio de } f
  \]

  \pause
  O que vai acontecer com outas retas da forma \(y=ax\)?
\end{resolution}

%------------------------------------------------------------------------------%
\begin{resolution}{Solução}
  Para qualquer $\delta>0$
  \pause
  sempre vai existir um

  ponto da forma \((x, x)\) e um ponto da forma \((x, 0)\)

  \pause
  dentro da bola de raio $\delta$ de centro em \((0, 0)\)

  \pause
  Nenhum valor $L$ atende as condições da definição de limite

  \pause
  \vspace{\baselineskip}
  \emph{O limite não existe}
\end{resolution}

%------------------------------------------------------------------------------%


%------------------------------------------------------------------------------%
\section{Inexistência do Limite}

%------------------------------------------------------------------------------%
\begin{frame}{Teste dos Dois Caminhos}
  Se \(f(x, y)\) tem limites diferentes ao longo de dois caminhos
  diferentes,

  no domínio de $f$,
  \pause
  quando $(x, y)$ se aproxima de $(x_0, y_0)$,

  \pause
  então
  \[
    \lim_{(x, y)\to(x_0, y_0)} f(x, y) \quad \text{ \emph{não existe}}
  \]
\end{frame}

%------------------------------------------------------------------------------%
\section{Exemplos}

%------------------------------------------------------------------------------%
\begin{question}
  Verifique se a função
  \(
  f(x, y) = \frac{2x^2y}{x^4+y^2}
  \)
  possui limite na origem.

  \pause
  \vspace{\baselineskip}
  O limite na origem leva a uma indeterminação do tipo $\sfrac{0}{0}$
\end{question}

%------------------------------------------------------------------------------%
\begin{resolution}{Solução}
  Fazendo \(y=mx\), para \(m\neq 0\) e \(x\neq 0\),
  \pause
  temos
  \[
    f(x, y)\evalat{y=mx}{}
    \pause
    \;=\; \frac{2x^2y}{x^4+y^2} \evalat{y=mx}{}
    \pause
    \;=\; \frac{2mx^3}{x^4+m^2x^2}
    \pause
    \;=\; \frac{2mx\,x^2}{\left(x^2+m^2\right)x^2}
    \pause
    \;=\; \frac{2mx}{x^2+m^2}
  \]

  \pause
  Calculando o limite \(x\to 0\), que corresponde a \((x, y)\to (0, 0)\)
  \[
    \pause
    \lim_{x\to 0} \frac{2mx}{x^2+m^2}
    \pause
    \;=\; \frac{0}{m^2}
    \pause
    \;=\; 0
    \qquad\qquad
    \forall m
  \]

  \pause
  Portanto,
  \(
    {
      \color<12->{red}
      \lim_{(x, y)\to(0, 0)} \frac{2x^2y}{x^4+y^2} = 0
    }
    \pause
    \qquad
    \qquad
    \text{\textcolor{red}{\textbf{Falso!}}}
  \)
\end{resolution}

%------------------------------------------------------------------------------%
\begin{resolution}{Solução}
  Fazendo \(y=kx^2\), para \(k\neq 0\) e \(x\neq 0\),
  \pause
  temos
  \[
    f(x, y)\evalat{y=kx^2}{}
    \pause
    \;=\; \frac{2x^2y}{x^4+y^2} \evalat{y=kx^2}{}
    \pause
    \;=\; \frac{2kx^4}{x^4+k^2x^4}
    \pause
    \;=\; \frac{2kx^4}{\left(1+k^2\right)x^4}
    \pause
    \;=\; \frac{2k}{1+k^2}
  \]

  \pause
  Calculando os limites até o ponto $(0, 0)$
  \[
    \pause
    \lim_{x\to 0} f(x, y)\evalat{y=kx^2}{}
    \pause
    \;=\; \lim_{x\to 0} \frac{2k}{1+k^2}
    \pause
    \;=\; \frac{2k}{1+k^2}
  \]
\end{resolution}

%------------------------------------------------------------------------------%
\begin{resolution}{Solução}
  Para cada valor de $k$, a função $f$ se aproxima da origem com um valor diferente

  \pause
  Portanto, o limite não existe
\end{resolution}

%------------------------------------------------------------------------------%

%------------------------------------------------------------------------------%
\begin{question}
  Verifique se a função
  \(
  f(x, y) = \frac{2x^2y}{x^4+y^2}
  \)
  possui limite na origem.

  \pause
  \vspace{\baselineskip}
  O limite na origem leva a uma indeterminação do tipo $\sfrac{0}{0}$
\end{question}

%------------------------------------------------------------------------------%
\begin{resolution}{Solução}
  Fazendo \(y=mx\), para \(m\neq 0\) e \(x\neq 0\),
  \pause
  temos
  \[
    f(x, mx)
    \pause
    \;=\; \frac{2x^2y}{x^4+y^2} \evalat{y=mx}{}
    \pause
    =\; \frac{2mx^3}{x^4+m^2x^2}
    \pause
    \;=\; \frac{2mx\,x^2}{\left(x^2+m^2\right)x^2}
    \pause
    \;=\; \frac{2mx}{x^2+m^2}
  \]

  \pause
  Calculando o limite \(x\to 0\), que corresponde a \((x, y)\to (0, 0)\)
  \[
    \pause
    \lim_{x\to 0} \frac{2mx}{x^2+m^2}
    \pause
    \;=\; \frac{0}{m^2}
    \pause
    \;=\; 0
    \qquad\qquad
    \forall m \neq 0
  \]

  \pause
  Portanto,
  \(
    {
      \color<12->{red}
      \lim_{(x, y)\to(0, 0)} \frac{2x^2y}{x^4+y^2} = 0
    }
    \pause
    \qquad
    \qquad
    \text{\textcolor{red}{\textbf{Falso!}}}
  \)
\end{resolution}

%------------------------------------------------------------------------------%
\begin{resolution}{Solução}
  Fazendo \(y=kx^2\), para \(x\neq 0\),
  \pause
  temos
  \[
    f(x, kx^2)
    \pause
    \;=\; \frac{2x^2y}{x^4+y^2} \evalat{y=kx^2}{}
    \pause
    =\; \frac{2kx^4}{x^4+k^2x^4}
    \pause
    \;=\; \frac{2kx^4}{\left(1+k^2\right)x^4}
    \pause
    \;=\; \frac{2k}{1+k^2}
  \]

  \pause
  Calculando os limites até o ponto $(0, 0)$
  \[
    \pause
    \lim_{x\to 0} f(x, y)\evalat{y=kx^2}{}
    \pause
    =\; \lim_{x\to 0} \frac{2k}{1+k^2}
    \pause
    \;=\; \frac{2k}{1+k^2}
  \]

  \pause
  Para cada valor de $k$, a função $f$ se aproxima da origem com um valor diferente

  \pause
  Portanto, o limite não existe
\end{resolution}

%------------------------------------------------------------------------------%

%------------------------------------------------------------------------------%
\begin{question}
  Mostre que o limite não existe
  \[
    \lim_{(x,y) \to (0, 0)} \frac{xy^2}{x^2 + y^4}
  \]
\end{question}

%------------------------------------------------------------------------------%
\begin{resolution}{Solução}
  Caminho \(x = y^2\)
  \[
    \pause
    \lim_{(x,y) \to (0,0)} \frac{xy^2}{x^2 + y^4} \evalat{x = y^2}{}
    \pause
    =\; \lim_{y \to 0} \frac{y^2 y^2}{y^4 + y^4}
    \pause
    \;=\; \lim_{y \to 0} \frac{y^4}{2y^4}
    \pause
    \;=\; \lim_{y \to 0} \frac{1}{2}
    \pause
    \;=\; \frac{1}{2}
  \]

  \pause
  Caminho \(x = 0\)
  \[
    \pause
    \lim_{(x,y) \to (0,0)} \frac{xy^2}{x^2 + y^4} \evalat{x = 0}{}
    \pause
    =\; \lim_{y \to 0} \frac{0 \times y^2}{0 + y^4}
    \pause
    \;=\; \lim_{y \to 0} 0
    \pause
    \;=\; 0
  \]

  \pause
  Como os limites nos caminhos são diferentes,
  \pause
  \emph{o limite não existe}
\end{resolution}

%------------------------------------------------------------------------------%

%------------------------------------------------------------------------------%
\begin{question}
  Mostre que o limite
  \(
    \lim_{(x,y) \to (0,0)} \frac{xy^2}{x^2 + y^4}
  \)
  não existe
\end{question}

%------------------------------------------------------------------------------%
\begin{resolution}{Solução}
  Caminho \(x = y^2\)
  \[
    \pause
    L_1
    \pause
    \;=\; \lim_{(x,y) \to (0,0)} \frac{xy^2}{x^2 + y^4} \evalat{x = y^2}{}
    \pause
    =\; \lim_{y \to 0} \frac{y^2 y^2}{y^4 + y^4}
    \pause
    \;=\; \lim_{y \to 0} \frac{y^4}{2y^4}
    \pause
    \;=\; \lim_{y \to 0} \frac{1}{2}
    \pause
    \;=\; \frac{1}{2}
  \]

  \pause
  Caminho \(x = 0\)
  \[
    \pause
    L_2
    \pause
    \;=\; \lim_{(x,y) \to (0,0)} \frac{xy^2}{x^2 + y^4} \evalat{x = 0}{}
    \pause
    =\; \lim_{y \to 0} \frac{0 \times y^2}{0 + y^4}
    \pause
    \;=\; \lim_{y \to 0} \frac{0}{y^4}
    \pause
    \;=\; \lim_{y \to 0} 0
    \pause
    \;=\; 0
  \]

  \pause
  Como \(L_1\neq L_2\),
  \pause
  o limite
  \(
    \lim_{(x,y) \to (0,0)} \frac{xy^2}{x^2 + y^4}
  \)
  \emph{não existe}
\end{resolution}

%------------------------------------------------------------------------------%


%------------------------------------------------------------------------------%
\section{Lista Mínima}

%------------------------------------------------------------------------------%
\begin{frame}{Lista Mínima}
  Cálculo Vol. 2 do Thomas 12\textsuperscript{a} ed. -- Seção 14.2
  \begin{enumerate}
    \item Estudar o texto da seção
    \item Resolver os exercícios: 41, 43, 45, 47
  \end{enumerate}

  \vfill
  Atenção: A prova é baseada no livro, não nas apresentações
\end{frame}

%------------------------------------------------------------------------------%
\end{document}
%------------------------------------------------------------------------------%
